% !PRINTED 535
\weberpage{535}
%<<CONT
(where now $A$, $B$ are of course not to be confused with the
coefficients of the quadratic form), then $B$ is an integral algebraic
number, and by § 74, (14) $A$ is the root of an algebraic equation whose
coefficients are integral algebraic numbers. Hence $A$ too is an integral
algebraic number.

But since the two quotients $u^r:v^r$ and $v^r:u^r$ are the roots of the
quadratic equation
\[
x^2 - Ax + 1 = 0
\]
it follows that
\[
\frac{u}{v} \quad \text{and} \quad \frac{v}{u}
\]
are integers, and, since they are reciprocal to one another, \gesperrt{units.}

\gesperrt{Thus $u$ and $v$ are associated numbers.}

Since one can now, by § 118, pass by repeated compositions with classes
$l$ (by which primes are representable) from every class $k$ to every other
class $k'$ of the same discriminant, and since two numbers associated with
a third are associated with one another, we have the theorem:

\gesperrt{If one puts for $\omega$ the $h$ roots of the forms of a system
of representatives of the discriminant $-4m$, then the $h$ numbers
$f(\omega)$ are associated with one another}; and from this, by (1), at
once the remarkable theorem, which may easily be confirmed in all the
examples:

\gesperrt{$f(\omega):2^\tau$ is an integer, and indeed a unit.}

1. The determination of the exponents $\tau$ is possible by
elementary means, if $m \equiv 1$ or $\equiv 2\ (\mathrm{mod}\ 4)$
or $m \equiv 3\ (\mathrm{mod}\ 8)$.

If in the equation with odd outer coefficients
\weq{3}{A\omega^2 + 2B\omega + C = 0}
with the discriminant $4(B^2 - AC) = -4m$ we make the substitution
\weq{4}{\omega = \frac{\omega'-1}{\omega'+1}, \quad \omega' = -\frac{\omega+1}{\omega-1},}
we obtain the equation
\weq{5}{\frac{A+2B+C}{2}\,\omega'^2 - (A-C)\,\omega' + \frac{A-2B+C}{2} = 0,}
%CONT>>
